% Options for packages loaded elsewhere
\PassOptionsToPackage{unicode}{hyperref}
\PassOptionsToPackage{hyphens}{url}
\PassOptionsToPackage{table}{xcolor}
%
\documentclass[
  letterpaper,
]{article}
\usepackage{amsmath,amssymb}
\usepackage{iftex}
\usepackage[T1]{fontenc}
\usepackage[utf8]{inputenc}
\usepackage{textcomp}
\usepackage{lmodern}
\ifPDFTeX\else
  % xetex/luatex font selection
\fi
% Use upquote if available, for straight quotes in verbatim environments
\IfFileExists{upquote.sty}{\usepackage{upquote}}{}
\IfFileExists{microtype.sty}{% use microtype if available
  \usepackage[]{microtype}
  \UseMicrotypeSet[protrusion]{basicmath} % disable protrusion for tt fonts
}{}
\makeatletter
\@ifundefined{KOMAClassName}{% if non-KOMA class
  \IfFileExists{parskip.sty}{%
    \usepackage{parskip}
  }{% else
    \setlength{\parindent}{0pt}
    \setlength{\parskip}{6pt plus 2pt minus 1pt}}
}{% if KOMA class
  \KOMAoptions{parskip=half}}
\makeatother
\usepackage{xcolor}
\usepackage{longtable,booktabs,array}
\usepackage{calc} % for calculating minipage widths
% Correct order of tables after \paragraph or \subparagraph
\usepackage{etoolbox}
\makeatletter
\patchcmd\longtable{\par}{\if@noskipsec\mbox{}\fi\par}{}{}
\makeatother
% Allow footnotes in longtable head/foot
\IfFileExists{footnotehyper.sty}{\usepackage{footnotehyper}}{\usepackage{footnote}}
\makesavenoteenv{longtable}
\setlength{\emergencystretch}{3em} % prevent overfull lines
\providecommand{\tightlist}{%
  \setlength{\itemsep}{0pt}\setlength{\parskip}{0pt}}
\setcounter{secnumdepth}{5}
% Self-contained presentation rules embedded in every generated LaTeX file.
\usepackage{iftex}
\ifPDFTeX\else\PackageError{SSE-ASAS}{Compile with pdfLaTeX only}{}\fi
\usepackage[T1]{fontenc}
\usepackage[utf8]{inputenc}
\usepackage{lmodern}
\usepackage{microtype}
\usepackage{geometry}
\geometry{margin=0.93in,headheight=16pt}
\usepackage{xcolor} % [table] passed at top of file (option-clash fix)
\definecolor{ink}{HTML}{15283B}
\definecolor{brand}{HTML}{075B75}
\definecolor{brandlight}{HTML}{EAF5F8}
\definecolor{warm}{HTML}{A15C1D}
\definecolor{muted}{HTML}{536577}
\usepackage{fancyhdr}
\usepackage{titlesec}
\usepackage{enumitem}
\usepackage{booktabs}
\usepackage{longtable}
\usepackage{array}
\usepackage{tabularx}
\usepackage{tcolorbox}
\tcbuselibrary{breakable,skins}
\usepackage{listings}
\usepackage{fancyvrb}
\usepackage{hyperref}
\hypersetup{colorlinks=true,linkcolor=brand,urlcolor=brand,citecolor=brand}
\urlstyle{same}
\usepackage{lastpage}
\setlength{\emergencystretch}{3em}
\setlength{\parskip}{0.45em}
\setlength{\parindent}{0pt}
\setlist[itemize]{leftmargin=*,topsep=3pt,itemsep=2pt}
\setlist[enumerate]{leftmargin=*,topsep=3pt,itemsep=2pt}
\titleformat{\section}{\Large\bfseries\color{ink}}{\thesection}{0.75em}{}[\vspace{-2pt}\color{brand}\titlerule]
\titleformat{\subsection}{\large\bfseries\color{brand}}{\thesubsection}{0.7em}{}
\titleformat{\subsubsection}{\normalsize\bfseries\color{ink}}{\thesubsubsection}{0.7em}{}
\pagestyle{fancy}
\fancyhf{}
\fancyhead[L]{\small\sffamily\textcolor{brand}{SOFTWARE SECURITY ENGINEERING}}
\fancyhead[R]{\small\sffamily\textcolor{muted}{TECHNICAL SERIES}}
\fancyfoot[L]{\scriptsize\sffamily Version 1.0 \textbullet\ 14 October 2026}
\fancyfoot[R]{\scriptsize\sffamily Page \thepage\ of \pageref{LastPage}}
\renewcommand{\headrulewidth}{0.5pt}
\renewcommand{\footrulewidth}{0.2pt}
\setcounter{tocdepth}{2}
\setcounter{secnumdepth}{3}
\lstdefinestyle{sse}{basicstyle=\ttfamily\footnotesize,frame=single,framerule=0.4pt,rulecolor=\color{brand!40},backgroundcolor=\color{brandlight!45},numbers=left,numberstyle=\scriptsize\color{muted},numbersep=6pt,breaklines=true,breakatwhitespace=false,tabsize=2,showstringspaces=false,columns=fullflexible,keepspaces=true,captionpos=b,aboveskip=9pt,belowskip=9pt}
\lstset{style=sse}
\newtcolorbox{sourceconcept}[1][]{enhanced,breakable,colback=brandlight,colframe=brand!65!black,title={ORIGINAL CONCEPT},fonttitle=\bfseries\sffamily,#1}
\newtcolorbox{modernextension}[1][]{enhanced,breakable,colback=white,colframe=warm!85!black,title={MODERN EXTENSION},fonttitle=\bfseries\sffamily,#1}
\newtcolorbox{fieldwork}[1][]{enhanced,breakable,colback=white,colframe=ink!55,title={PRACTICAL APPLICATION},fonttitle=\bfseries\sffamily,#1}
\AtBeginDocument{\color{ink}}
\makeatletter
\renewcommand{\maketitle}{%
 \begin{titlepage}\thispagestyle{empty}
 \vspace*{1.5cm}{\color{brand}\rule{\textwidth}{3pt}}\par\vspace{1cm}
 {\sffamily\bfseries\color{muted}\large SOFTWARE SECURITY ENGINEERING\\ APPLICATION SECURITY ASSESSMENT SERIES}\par\vspace{1.05cm}
 {\sffamily\bfseries\color{ink}\fontsize{27}{34}\selectfont\@title\par}\vspace{0.8cm}
 {\sffamily\large\color{brand}\@author\par}\vfill
 \begin{tcolorbox}[colback=brandlight,colframe=brand!55,boxrule=0.5pt]
 \small\textbf{Editorial policy}\quad Original teaching synthesis with book-grounded concepts distinguished from contemporary practices. Modern platform behavior is version-sensitive; validate against linked vendor documentation before production use.
 \end{tcolorbox}
 {\color{brand}\rule{\textwidth}{2pt}}\par\vspace{0.3cm}
 {\sffamily\small Version 1.0 \quad|\quad October 2026 \quad|\quad pdfLaTeX}\end{titlepage}\setcounter{page}{1}}
\makeatother
\ifLuaTeX
  \usepackage{selnolig}  % disable illegal ligatures
\fi
\usepackage{bookmark}
\IfFileExists{xurl.sty}{\usepackage{xurl}}{} % add URL line breaks if available
\urlstyle{same}
\hypersetup{
  pdftitle={Security Foundations and Assurance},
  pdfauthor={Independent Technical Study Edition},
  hidelinks,
  pdfcreator={LaTeX via pandoc}}

\title{Security Foundations and Assurance}
\author{Independent Technical Study Edition}
\date{}

\begin{document}
\maketitle

{
\setcounter{tocdepth}{3}
\tableofcontents
}
\section{Purpose, competence, and reading
route}\label{purpose-competence-and-reading-route}

This guide establishes the formal vocabulary for the rest of the series.
An engineer completing it should be able to state security properties,
model access decisions, distinguish policy from enforcement, identify
trust assumptions, and construct an assurance argument rather than
equating a passing scan with security.

\textbf{Baseline:} Matt Bishop, \emph{Computer Security: Art and
Science}, second edition, Chapters 1--9, 14--22, and 24--25. Chapters
2--3 develop protection-state models, Chapters 4--9 develop policies,
and Chapters 19--22 develop assurance. \textbf{Contemporary layer:}
application authorization, workload identity, distributed systems, and
NIST SSDF practices are applications of those principles, not material
attributed to the original text.

\section{Security as a property of a
system}\label{security-as-a-property-of-a-system}

\subsection{Assets, policy, and the protection
state}\label{assets-policy-and-the-protection-state}

Let \(S\) be subjects, \(O\) objects, and \(R\) rights. The \emph{access
control matrix} is a function \(A:S\times O\to 2^R\); for example
\(A(\text{analyst},\text{report})=\{\text{read}\}\). A protection state
specifies which operations are authorized at a particular instant. A
transition changes that state, for example by adding a right to a role
or by revoking a credential.

\begin{sourceconcept}
The key distinction is **policy** (which states may occur and which accesses are permitted), **mechanism** (the implementation enforcing the rule), and **assurance** (evidence that the implementation actually enforces the intended policy). A secure mechanism that enforces the wrong policy does not meet system requirements.
\end{sourceconcept}

Consider an invoice service in which an authenticated user can read an
invoice by numeric ID. Authentication answers \emph{who} is asking.
Authorization must also verify that the invoice belongs to the user's
tenant. A lookup that checks only whether the user is logged in violates
the application's protection policy even if the HTTP endpoint is well
authenticated.

\subsection{Confidentiality, integrity, and
availability}\label{confidentiality-integrity-and-availability}

\begin{itemize}
\tightlist
\item
  \textbf{Confidentiality:} unauthorized information disclosure does not
  occur. Examples include tenant isolation and protection of signing
  keys.
\item
  \textbf{Integrity:} unauthorized modification of data or processing is
  prevented or detected. Examples include restricting payment approval
  and maintaining an append-only audit trail.
\item
  \textbf{Availability:} legitimate requests receive needed service
  within agreed resource and timing bounds. Examples include admission
  control, resource quotas, and resilient deployment.
\end{itemize}

Security also requires authenticity, accountability, privacy, and
appropriate operation under failure. These are not automatically
entailed by the CIA triad. Define explicit expected outcomes rather than
relying on labels.

\begin{longtable}[]{@{}
  >{\raggedright\arraybackslash}p{(\columnwidth - 4\tabcolsep) * \real{0.3333}}
  >{\raggedright\arraybackslash}p{(\columnwidth - 4\tabcolsep) * \real{0.3333}}
  >{\raggedright\arraybackslash}p{(\columnwidth - 4\tabcolsep) * \real{0.3333}}@{}}
\toprule\noalign{}
\begin{minipage}[b]{\linewidth}\raggedright
Property
\end{minipage} & \begin{minipage}[b]{\linewidth}\raggedright
Measurable expectation
\end{minipage} & \begin{minipage}[b]{\linewidth}\raggedright
Adversarial question
\end{minipage} \\
\midrule\noalign{}
\endhead
\bottomrule\noalign{}
\endlastfoot
Confidentiality & User A cannot read tenant B data & Can a request
select an object outside its authorization scope? \\
Integrity & Only approved workflows alter settlement status & Can a user
bypass dual approval via an alternate API? \\
Availability & Service degrades predictably under overload & Can
expensive requests starve normal requests? \\
Accountability & Decisions have attributable, tamper-evident events &
Can actions occur without traceable subject identity? \\
\end{longtable}

\section{Models of permitted
behavior}\label{models-of-permitted-behavior}

\subsection{Discretionary and mandatory access
controls}\label{discretionary-and-mandatory-access-controls}

A discretionary model places control of some grants with resource
owners; a mandatory model constrains grants by a governing
classification or other global rule. A practical system can combine
discretionary membership with non-bypassable tenant separation. Lists of
permissions, capabilities, RBAC, and attribute-based checks are
mechanisms or policy styles with different delegation properties.

\textbf{Bell--LaPadula} models confidentiality using security
classifications and information-flow constraints. Its familiar
simplified intuitions are \emph{no read up} and \emph{no write down}
under the corresponding conditions. These are statements about a
particular confidentiality policy; they do not prove integrity,
availability, or application correctness.

\textbf{Biba} concerns integrity levels, with constraints that oppose
contamination of high-integrity computations by lower-integrity input in
its strict form. \textbf{Clark--Wilson} treats integrity through
well-formed transactions, constrained data items, verification, and
separation of duties; it is particularly instructive for business
workflows. The \textbf{Chinese Wall} model handles conflict-of-interest
policies in which access history influences subsequent access. These
models cannot be treated as interchangeable.

\begin{longtable}[]{@{}
  >{\raggedright\arraybackslash}p{(\columnwidth - 4\tabcolsep) * \real{0.3333}}
  >{\raggedright\arraybackslash}p{(\columnwidth - 4\tabcolsep) * \real{0.3333}}
  >{\raggedright\arraybackslash}p{(\columnwidth - 4\tabcolsep) * \real{0.3333}}@{}}
\toprule\noalign{}
\begin{minipage}[b]{\linewidth}\raggedright
Model
\end{minipage} & \begin{minipage}[b]{\linewidth}\raggedright
Central question
\end{minipage} & \begin{minipage}[b]{\linewidth}\raggedright
Application security exercise
\end{minipage} \\
\midrule\noalign{}
\endhead
\bottomrule\noalign{}
\endlastfoot
Access matrix & Which right does each subject have on each object? &
Build an endpoint-by-role decision matrix \\
Bell--LaPadula & May information cross a classification boundary? &
Trace export and reporting data flows \\
Biba & Could untrusted input contaminate a trusted decision? & Protect a
release-signing pipeline \\
Clark--Wilson & Does the action preserve a valid business state? &
Enforce review and approval of payment changes \\
Chinese Wall & Does historical access create a conflict? & Restrict
consultants to one competing client data set \\
\end{longtable}

\subsection{Reference monitor and complete
mediation}\label{reference-monitor-and-complete-mediation}

A reference monitor is the conceptual intermediary for subject/object
access. Its enforcement implementation must be tamper resistant, invoked
on every relevant access, and small enough to evaluate. This becomes a
powerful design test: where does an authorization check live, which
paths bypass it, and what evidence demonstrates it executes before every
sensitive action?

\begin{modernextension}
In a microservice architecture, a gateway may authenticate a request, but the resource-owning service still needs object-level authorization. Authorization claims forwarded over untrusted hops are not enforcement by themselves. Short-lived service identity, service-to-service authentication, deny-by-default policy, and request-context propagation are modern engineering applications of complete mediation.
\end{modernextension}

\subsection{Information flow, covert channels, and
composition}\label{information-flow-covert-channels-and-composition}

A function-level access decision is narrower than an information-flow
guarantee. A service may deny \texttt{GET\ /report} while leaking report
content through verbose errors, logs, caches, metrics, search
autocomplete, or timing. Composition of locally correct components can
violate a system-wide property if assumptions differ across trust
boundaries.

A useful invariant is: \emph{no response field, log field, or outbound
event contains data classified above the requesting principal's
authorized level}. The audit scope must include indirect paths. State
assumptions explicitly, such as trusted clock, secure key store, correct
tenant mapping, and confidentiality of system logs.

\section{Threats, assumptions, and
assurance}\label{threats-assumptions-and-assurance}

\subsection{From security expectation to threat
model}\label{from-security-expectation-to-threat-model}

\begin{enumerate}
\def\labelenumi{\arabic{enumi}.}
\tightlist
\item
  Declare the protected asset and expected invariant.
\item
  Identify subjects, entry points, data stores, processes, and trust
  boundaries.
\item
  Enumerate adverse state transitions and unauthorized information
  flows.
\item
  Record attacker-controlled inputs, capabilities, preconditions, and
  impact.
\item
  Select preventive and detective mechanisms that address those
  transitions.
\item
  Specify testable evidence and residual assumptions.
\end{enumerate}

A threat is not identical to a vulnerability: the former describes a
potential adverse event or actor capability; the latter is a weakness
that may allow a threat to succeed. An incident is observed occurrence,
and risk incorporates likelihood, consequence, and uncertainty.

\subsection{Assurance argument}\label{assurance-argument}

A persuasive claim is supported by multiple kinds of evidence. For a
tenant-isolation claim, evidence might include a documented policy, a
centralized ownership guard, negative unit tests for cross-tenant
access, integration tests across endpoints, access logs showing denied
attempts, and a review of bypass paths. A scanner's absence of findings
cannot substitute for these.

\begin{longtable}[]{@{}
  >{\raggedright\arraybackslash}p{(\columnwidth - 4\tabcolsep) * \real{0.3333}}
  >{\raggedright\arraybackslash}p{(\columnwidth - 4\tabcolsep) * \real{0.3333}}
  >{\raggedright\arraybackslash}p{(\columnwidth - 4\tabcolsep) * \real{0.3333}}@{}}
\toprule\noalign{}
\begin{minipage}[b]{\linewidth}\raggedright
Assurance claim
\end{minipage} & \begin{minipage}[b]{\linewidth}\raggedright
Evidence
\end{minipage} & \begin{minipage}[b]{\linewidth}\raggedright
Limit or counterexample
\end{minipage} \\
\midrule\noalign{}
\endhead
\bottomrule\noalign{}
\endlastfoot
Every invoice read is authorized & Route inventory; authorization
middleware; negative tests & Direct database workers may bypass HTTP
layer \\
Approval is atomic & Transaction tests; concurrency tests; data
constraints & Side effects in queues may not share transaction \\
Security policy is reviewed & Versioned requirements; signed design
decision & Stale policy after a new API ships \\
\end{longtable}

\begin{sourceconcept}
Bishop's treatment of assurance explicitly separates specification, design, implementation, evaluation, and the trusted computing base. The size and complexity of the enforcement boundary affect how much confidence analysis and testing can justify.
\end{sourceconcept}

\section{Applied model: multi-tenant invoice
platform}\label{applied-model-multi-tenant-invoice-platform}

Define \(T(u)\) as the set of tenants that principal \(u\) may access,
and \(tenant(o)\) as the tenant owning object \(o\). For a read request,
use

\[allow(u,o,\mathrm{read}) \Longleftrightarrow authenticated(u)\ \land\ tenant(o)\in T(u)\ \land\ permission(u,\mathrm{readInvoice}).\]

For a payment approval, add a workflow invariant: the submitter cannot
be the final approver and the invoice must be in an approvable state.
These constraints are separate from read authorization. Compare them
with Bell--LaPadula's confidentiality concern and Clark--Wilson's
transactional integrity concern; the analogy illuminates design, but it
is not a claim that either formal model directly specifies a SaaS
product.

\begin{fieldwork}
**Design review task.** List all ways an invoice can leave the platform: API reads, PDF download, bulk export, notification, message queue, search index, logging, support tooling, backups, and analytics. For each path, name the enforcement point and identify whether the claim is directly tested. Record unresolved trust assumptions.
\end{fieldwork}

\section{Engineering controls and success
criteria}\label{engineering-controls-and-success-criteria}

\begin{itemize}
\tightlist
\item
  Store authorization decisions close to protected resources; fail
  closed on missing context.
\item
  Normalize identities across human, workload, automation, and vendor
  principals.
\item
  Re-evaluate authorization on every sensitive operation; avoid relying
  on UI hiding.
\item
  Make policy test cases explicit and version-controlled, including
  negative tests.
\item
  Separate security-sensitive responsibilities and require dual approval
  where justified.
\item
  Restrict secrets, logs, caches, backups, and telemetry under the same
  information-classification policy.
\item
  Validate policy during design and after change; record both residual
  risks and verification limits.
\end{itemize}

\textbf{Exit criteria:} deliver a one-page policy statement, access
matrix, two data-flow diagrams with trust boundaries, a set of at least
ten negative authorization test cases, and a structured assurance
argument linking each claim to evidence.

\section{Annex A: policy engineering
workbook}\label{annex-a-policy-engineering-workbook}

\subsection{Turn stakeholder statements into enforceable
invariants}\label{turn-stakeholder-statements-into-enforceable-invariants}

Stakeholder statement: ``Customer documents are private.'' This is not
yet a security policy; it does not identify owner, sharing, metadata,
support privileges, derivative products, or deletion behavior. A
testable refinement is: \emph{For every object \(o\), if
\(tenant(o)\neq tenant(u)\) and principal \(u\) has no explicitly
reviewed cross-tenant privilege, read, list, export, search-index,
background-generation, and metadata APIs shall not expose \(o\) or any
data derived from \(o\).} The word ``derived'' prevents the policy from
omitting thumbnail previews, audit event payloads, and full-text search
results.

Create a policy worksheet with these columns: controlled resource,
requesting principal, required action, allowed preconditions, denied
conditions, enforcement implementation, tests, and residual trust
assumptions. The same worksheet doubles as an implementation review and
as an evidence-based assurance claim. A good policy avoids dependence on
a single endpoint name so it can survive API refactoring.

\subsection{RBAC, attributes, and delegated
authority}\label{rbac-attributes-and-delegated-authority}

RBAC is appropriate when permissions naturally follow job functions and
stable roles. Attribute-based decisions can refine them using tenant,
object ownership, workflow status, classification, device or service
identity, and authorization context. The roles alone must never be
interpreted as proof of object ownership. A ``support'' role is
especially sensitive: require explicit support scope, audit trail,
expiration, and justification for cross-tenant access. In some
architectures, a capability token can delegate authority; constrain its
resource, operation, audience, expiration, and revocability.

A policy decision can be decomposed as
\(allow = identityValid \land rolePermits \land objectMatches \land statePermits \land contextualConstraints\).
This Boolean illustration assumes conjunctive policy. Real systems may
use explicitly reviewed exceptions; never allow a catch-all
administrator clause to unintentionally override tenant-isolation
invariants.

\subsection{A claim-evidence-risk
case}\label{a-claim-evidence-risk-case}

\textbf{Top-level claim C1:} No user can access another tenant's
confidential invoice without an authorized support workflow.
\textbf{Subclaim C1.1:} The service authenticates the user and binds the
tenant to the server-issued principal. \textbf{Subclaim C1.2:} Every
invoice read and export checks tenant ownership at the authoritative
query. \textbf{Subclaim C1.3:} Support accesses require temporary
approval and protected audit records.

\textbf{Evidence for C1.1:} identity middleware configuration, token
validation tests, and a forged-tenant negative test. \textbf{For C1.2:}
endpoint inventory, repository search of invoice query APIs, tests for
read/list/export and background jobs, and a design assertion that direct
SQL access is limited. \textbf{For C1.3:} policy configuration, approval
workflow test, revocation test, and a log integrity review.

\textbf{Defeaters and uncertainty:} unreviewed internal workers,
eventual consistency in caches, external analytics feeds, logs that
disclose customer data, and policy changes outside version control.
Treat each defeater as a review task, explicit limitation, or residual
risk, not as quietly satisfied.

\subsection{Privacy, availability, and usability
tensions}\label{privacy-availability-and-usability-tensions}

Confidentiality controls can impose availability constraints; for
instance, a key management outage may prevent decryption of critical
records. An authorization service can become a single point of failure.
Specify fail-closed behavior together with resilient policy caching,
short bounded decision freshness, emergency access controls, and
observability. Emergency access should be narrowly scoped, attributable,
time-limited, and reviewed after use, rather than a permanent bypass.

\textbf{Workbook checks:} at least one negative test for every
high-impact permission; tests for stale attributes and revoked access;
policy-change review; protected logs; secure behavior during missing
authentication context; performance/failure budget for policy checks; a
documented risk owner for each unverified trust assumption.

\section{Source and modern extension
map}\label{source-and-modern-extension-map}

\textbf{Book-grounded:} Bishop, Chapters 1--3 (components and protection
models), 4--9 (policies and information flow), 14--18 (design
principles, identity, enforcement, confinement), 19--22 (assurance), and
24--25 (vulnerability analysis and auditing). Read chapter sections for
the complete definitions, variants, and proofs.

\textbf{Modern extension:} NIST SP 800-218 v1.1:
\url{https://csrc.nist.gov/pubs/sp/800/218/final}; OWASP ASVS:
\url{https://owasp.org/projects/asvs}; OWASP SAMM:
\url{https://owasp.org/projects/samm}. These resources augment rather
than replace the source models.

\end{document}
